\documentclass[11pt,a4paper]{article}
\usepackage[T1]{fontenc}
\usepackage{lmodern,amsmath,amssymb,amsthm,booktabs,microtype}
\usepackage[margin=27mm]{geometry}
\usepackage[colorlinks=true,linkcolor=blue,citecolor=blue,urlcolor=blue]{hyperref}
\setlength{\emergencystretch}{3em}
\newtheorem{theorem}{Theorem}[section]
\newtheorem{lemma}[theorem]{Lemma}
\newtheorem{proposition}[theorem]{Proposition}
\newcommand{\Q}{\mathbb Q}\newcommand{\Z}{\mathbb Z}
\newcommand{\Tr}{\operatorname{Tr}}\newcommand{\Nm}{\operatorname{N}}
\newcommand{\ord}{\operatorname{ord}}\newcommand{\PP}{\mathbb P}
\title{Native pole witnesses: norm composition and forced ramification}
\author{}\date{17 November 2026}
\begin{document}\maketitle
\begin{abstract}
The detector product has an explicit quadratic norm representation.
An actual seven-contact multiplication tuple and physical slope show that
this representation does not supply a native pole witness: all pre-pole
whole-character guards hold, but the pole input has valuation $-7$ at
$100003$. We also show that the native pole equations themselves, over
discretely valued function fields, force even ramification above every
negative odd contact pole. A degree-$d$ fixed-point slope cover carrying
native witnesses therefore has genus at least $1+d(r-4)/4$, where
$r=10,14,18$ for the general, seven-contact and nine-contact sources.
The existing abelian-family argument extends to native pole witnesses,
including the even-valuation revisions. These conclusions concern
specific geometric constructions; they do not disprove prepared rational
realization or give a lower bound for other Diophantine definitions.
\end{abstract}

\section{What the built-in norm actually gives}
Let $K=\Q(\sqrt2)$, let $u\in\Q^\times$, and let $D_i\in K$.
In the quadratic algebra with $s^2=-u$, put
\[
 Z=\prod_i(1+D_i s/u)=X+Ys,
 \qquad X=a+b\sqrt2,\quad Y=c+d\sqrt2.
\]
Then
\begin{equation}\label{eq:norm}
 W=\prod_i(1+D_i^2/u)=X^2+uY^2=W_0+W_1\sqrt2,
\end{equation}
where
\[
 W_0=a^2+2b^2+u(c^2+2d^2),\qquad W_1=2(ab+ucd).
\]
Taking the other quadratic norm first gives
\begin{equation}\label{eq:cross}
 \Nm_{K/\Q}(W)=C^2+uE^2=W_0^2-2W_1^2,
 \quad C=a^2-2b^2-u(c^2-2d^2),\quad E=2(ac-2bd).
\end{equation}
Thus, whenever $W_0\ne0$, there is an explicit rational point on
\[
              (C/W_0)^2+u(E/W_0)^2+2(W_1/W_0)^2=1.
\]
Writing $Z^{-1}=A+Bs$ likewise expresses
$\Tr(W^{-1})$ as twice a value of
$x_1^2+2x_2^2+u(x_3^2+2x_4^2)$.
These identities hold for every list $D_i$. They add no condition to
the detector. The following example shows why they cannot be used as
an automatic native pole certificate, even on actual multiplication points.

\section{An actual physical counterexample}
Use the original squared balancing of the seven-contact source
\cite{Seven}, with its enlarged auxiliary charts. On
\[
 E:y^2=x^3-(11-6\sqrt2)x
\]
take
\[
 X=1,\quad B=3,\quad A=P,\quad H=U=R=(1+\sqrt2,2\sqrt2),
 \quad 2R=P.
\]
This is the intended multiplication relation $U=[1]H$.
Its selected-character tests are
\[
 (C_1,\ldots,C_6)=(c,c,\bar c,\bar c,0,0),\qquad c=\sqrt2-1.
\]
Set
\begin{equation}\label{eq:fixture}
 p=100003,\quad M=\frac{10621537}{256},\quad a=M^8,
 \quad z=a+8+p,
\end{equation}
and use the actual physical functions
\[
 F=\prod_{c_0\in\{-7,-6,-4,1,3,5,8\}}((a+c_0)^2-z^2),
 \quad u=-14F/F_a,\quad w=\frac{u}{(1+u)^2},
\]
where differentiation precedes specialization of $a$.
The number $p$ is prime. Exact arithmetic gives the residues below;
the modulus is $p^2=10000600009$.
\[
\begin{array}{c|rrrr}
 &a&z&F&F_a\\\hline
 \bmod p^2&4435335376&4435435387&3047091410&2746251916
\end{array}
\]
In particular $v_p(F)=1$, $v_p(F_a)=0$, $v_p(u)=1$, and
$(u/p)\bmod p=14$. Also $a>10$ and $z>a+10$.
Every factor of $F$ is negative and every summand in $F_a$ is
positive. More precisely,
\[
 u=\frac{14}{\sum_{c_0}2(a+c_0)/(z^2-(a+c_0)^2)}>2p>0,
\]
because each denominator difference $z-a-c_0$ is at least $p$.

For this actual tuple define
\[
 D(u)=u^2(u^2+6u+1)^2,\qquad
 R(u)=D(u)+2u^3+12u^2+34u+1.
\]
Direct conjugate multiplication, independently reconstructed in
\cite{Seven}, gives
\[
 W=(R/D)^2,\qquad V=\Tr(W^{-1})=2/W,\qquad
                       b=\frac{4w}{V^2}=wW^2.
\]
Consequently
\begin{equation}\label{eq:bad}
 v_p(D)=2,\quad v_p(R)=0,\quad v_p(W)=-4,
 \quad v_p(V)=4,\quad v_p(b)=v_p(7b)=-7.
\end{equation}
The norm product \eqref{eq:norm} still holds exactly: its factors use
$D_i=C_i^2/(1+uC_i^2)$ with the actual six tests above.

\begin{proposition}\label{prop:example}
The tuple \eqref{eq:fixture} satisfies every pre-pole point, contact,
balanced-denominator, product and trace guard of the seven-contact
source at all $64$ geometric point characters. Its built-in norm
representation does not lift to any intrinsic native pole witness.
\end{proposition}
\begin{proof}
At each real slot the three point ordinates have independent signs.
Every required difference is one of $\pm R,\pm3R$. Exact duplication
and addition verify that $P,R,3R$ have nonzero abscissae and ordinates
at both embeddings, and $x(P)-x(R)$ is nonzero at both. The old point,
difference and quotient guards therefore hold at every character.
The scale and contact functions depend only on retained abscissae,
so the positivity just proved holds at every character.

All balanced denominators are $1+uC_i^2>0$ at every real slot.
Every detector factor is at least one; thus products, their norms,
and their reciprocal traces are positive. The trace and intrinsic
input are nonzero. The companion checker also evaluates the actual
stored coefficient DAG at all $64$ characters modulo $1000000007$.
Every native pre-pole guard has nonzero reduction, independently
certifying its nonvanishing in characteristic zero. The checker
verifies all $768$ balanced denominator signs exactly in $\Q(\sqrt2)$.

Finally \eqref{eq:bad} and Lemma~\ref{lem:intrinsic} below exclude
every intrinsic pole witness, including one with zero final ordinate.
\end{proof}

The example is not asserted to meet fixed-place or selected-prime
preparation. It shows failure of an automatic implication from the
norm identity, actual multiplication points, a physical slope, and
all pre-pole guards. It does not show failure of the integer input's
prepared existence assertion. The fixture $B=3$ is not certified to
be a sufficient defining parameter.

\section{Native pole equations over valued function fields}
All valuations in this section are discrete, with values in $\Z$.
The value of zero is $+\infty$.

\begin{lemma}\label{lem:intrinsic}
Over any discretely valued field, a solution of
\begin{equation}\label{eq:intrinsic}
 Y^2=\eta(\eta+v^2)^2s^4-
                 \eta(\eta-v^2)^2(1-\eta)^2
\end{equation}
cannot have a negative odd value of $\eta$.
No nonzero restriction on $v,s,Y$ is needed.
\end{lemma}
\begin{proof}
Suppose $\nu=v(\eta)<0$ is odd. Set
$r=\min(\nu,2v(v))$. The two terms defining $\eta\pm v^2$
have unequal valuations, so both expressions have value $r$.
Also $v(1-\eta)=\nu$. The two summands on the right of
\eqref{eq:intrinsic} have values
\[
                       \nu+2r+4v(s),\qquad3\nu+2r.
\]
When $s\ne0$ they are odd and unequal, since equality would give
$2v(s)=\nu$. When $s=0$, only the second remains. The radicand
is therefore nonzero with odd valuation, contradicting the square.
This is the proof of \cite[Theorem 1.1]{Intrinsic}, with no use of
the residue field being finite or of the valued field being $\Q_p$.
\end{proof}

\begin{lemma}\label{lem:general}
Suppose $v(2)=v(3)=v(17)=0$. The general native equations
\begin{equation}\label{eq:general}
 p_0^2=B_0(1-e)(1+3e),\qquad
 B_0q_0^2=(s^2-B_0e(B_0-b^2))(s^2+3B_0e(B_0-b^2)),
 \quad B_0=17b
\end{equation}
exclude a negative odd valuation of $b$. This permits zero ordinates.
The same assertion holds in their guarded rescaled chart.
\end{lemma}
\begin{proof}
Let $\nu=v(b)<0$ be odd. If $v(e)<0$, the first right side has
valuation $\nu+2v(e)$, which is odd. If $v(e)>0$, including
$e=0$, it has valuation $\nu$, again odd. Thus $v(e)=0$.
Write $A_0=B_0e(B_0-b^2)$. Since $v(B_0-b^2)=2\nu$, we get
$v(A_0)=3\nu$, an odd integer. The values $2v(s)$ and $3\nu$
are unequal. Both factors in the second right side have valuation
$\min(2v(s),3\nu)$; their product is nonzero of even valuation.
The left side is zero or has odd valuation, a contradiction.

The rational rescaling in \cite[Section 4]{Unit} is invertible when
its input, multiplier and denominator guards hold. Hence it carries
any rescaled-chart solution back to \eqref{eq:general}. The argument
does not require the last ordinate to be nonzero.
\end{proof}

For geometric valuations in characteristic zero, rational nonzero
constants have valuation zero. Thus both lemmas apply, and replacing
$b$ by $7b$ in the intrinsic source does not change a pole order.

\section{Every finite slope cover carrying native witnesses}

\begin{theorem}\label{thm:cover}
Let $k$ be an algebraically closed field of characteristic zero and
let $b(t)\in k(t)^\times$ have negative odd order at $r\ge1$ distinct
points of $\PP^1$. Let $C$ be a smooth projective connected curve
with a finite map $t:C\to\PP^1$ of degree $d$, and suppose
$k(C)$ contains native pole witnesses for $b(t)$ of either preceding
kind, on the required guarded chart. Then every ramification index
above these $r$ points is even. In particular $d$ is even and
\begin{equation}\label{eq:genus}
                         g(C)\ge1+\frac{d(r-4)}4.
\end{equation}
The same conclusion holds for any replacement certificate whose
solutions over these valued function fields exclude negative odd
orders of the same input.
\end{theorem}
\begin{proof}
At a point of $C$ above a pole of odd order $m<0$, the pulled-back
order is $em$, where $e$ is its ramification index. The native local
lemmas force $e$ even. Thus there are at most $d/2$ geometric points
in each such fiber, and its contribution to ramification is at least
$d-d/2=d/2$. Riemann--Hurwitz \cite{Stacks} gives
\[
                      2g(C)-2\ge-2d+rd/2,
\]
which is \eqref{eq:genus}. The fibers also show that $d$ is even.
\end{proof}

Fix a guarded rational point character in one of the actual detector
sources, including its point coordinates and scale. In a derivative
source $a>\max_c|c|$, as in the actual construction, so its $2k$
contact positions are pairwise distinct. At a derivative contact
$z=\pm(a+c)$, the derivative
weight has a simple zero: for $k$ shifts,
\[
 u=-2kF/F_a,\qquad u'(\pm(a+c))=\pm2k.
\]
The same is true of $w=u/(1+u)^2$. Every nonzero balanced test
contributes a simple pole to $W$ in each slot; each zero test
contributes one. The required reciprocal tests are nonzero, so
$V=\Tr(W^{-1})$ vanishes to order at least one. If cancellation
raises that order, the following conclusion only becomes stronger:
\[
                  \ord(b)=1-2\ord(V)<0\quad\hbox{is odd}.
\]
If $V$ vanishes identically, the guarded fiber is empty and is omitted.
The argument applies to both versions of balanced abscissa tests.

For the first general detector, its ten contacts are simple zeros
of $q=\rho R$ and of its normalized $w$. The normalizer is a unit.
Each nonzero field factor has a pole; zero field tests contribute
one. Scalar factors tend to one, with second summands of order at
least three in the native source and at least five in the normalized
source. The nonzero reference reciprocal supplies a pole in each
slot. Therefore the same negative odd order of $b$ follows. It also
holds for the unsquared first-ten revision \cite{Even25}.

The resulting bounds are
\[
\begin{array}{lrrl}
\toprule
\text{Source}&r&\text{minimum }g\ (d=2)&\text{bound for degree }d\\
\midrule
\text{General first detector}&10&4&1+3d/2\\
\text{Seven derivative contacts}&14&6&1+5d/2\\
\text{Nine derivative contacts}&18&8&1+7d/2\\
\bottomrule
\end{array}
\]
This excludes every genus-zero or genus-one finite slope cover
carrying native witnesses, regardless of which rational expressions
are used for its pole parameters. It is not restricted to the
previous actual-trace square equation. A genus bound does not exclude
individual rational points or higher-dimensional arithmetic families.

\section{Nonconstant slopes on abelian families}
Let $\mathcal R=\operatorname{Res}_{K/\Q}E$. We use the established
simplicity of $\mathcal R$ over $\Q$ and the point-map facts in
\cite[Sections 3 and 5]{Abelian}. They concern the original and full
auxiliary point domains. The following proof spells out the additional
native-pole step and the required codimension condition.

\begin{theorem}\label{thm:abelian}
Let $J/\Q$ be an abelian variety with a rational map into a guarded
native-pole source. In the derivative source fix the input, coefficient
parameter and root point $A$, so its contact positions are constant.
In the general source fix the input and coefficient exponent; its
first detector already has constant comparison scale one. Then the
physical slope of the rational map is constant. This includes both
even-valuation revisions, and holds after isogenies. No exclusion of
constant-slope abelian families is asserted.
\end{theorem}
\begin{proof}
The reference-point map $J\dashrightarrow\mathcal R$ extends to
a morphism, a homomorphism followed by translation
\cite[Chapter I, Theorem 3.2 and Corollary 3.7]{Milne}.
By simplicity its image is a point or all of $\mathcal R$.
If constant, the two reference reciprocals are nonzero constants by
the guards. Otherwise work over $\overline\Q$, where
$\mathcal R=E\times E^\sigma$. In the first slot the simultaneous
zero locus of
\[
                     f(H),\qquad\overline{f(H-P)}
\]
is the inverse image of $(O,\bar P)$; in the other slot it is the
inverse image of $(P,O)$. The only zero of $f=1/x$ is $O$.
Fibers of a surjective abelian homomorphism have codimension two.
Thus no divisor has both reference reciprocals of positive order
in either slot.

Consider a prime divisor of $J_{\overline\Q}$ where the slope
minus one of its constant contact values has positive odd order
$\nu$. The physical contact weight has order $\nu$.
In the derivative source, every detector factor has nonpositive
order. For the original balancing this follows from
\[
 \ord(D_i)=\max(2\ord(C_i),-\nu),\qquad
 \ord(1+D_i^2/u)=\min(0,2\ord(D_i)-\nu).
\]
For unsquared abscissa balancing replace $2\ord(C_i)$ in the first
formula by $\ord(C_i)$. This is valid because the divisor of an
elliptic reciprocal abscissa is even, also after pullback. In each
slot one reference test has nonpositive order, so its factor has
order at most $-\nu$.

For the general source the normalizer is a unit at a contact.
Its squared factors have an even numerator order minus an odd
multiple of $\nu$. The unsquared reciprocal factors have the same
parity because their reciprocal divisors are even. All factors again
have nonpositive order, and one reference factor in each slot has
order at most $-\nu$. Zero comparison or scalar tests contribute one.

In either source $\ord(V)\ge\nu$, including any cancellation
between slots. Hence the pole input has odd negative order
\[
                         \ord(b)=\nu-2\ord(V)\le-\nu.
\]
The native pole lemmas contradict the supposed rational witnesses.
Every component of each contact fiber must therefore have even
multiplicity.

Let $t$ denote the slope and let $t_1,\ldots,t_r$ be these constant
contacts. The divisor of $q=\prod_i(t-t_i)$ is even: its zero
divisors are even and its poles are $r$ times the pole divisor of
$t$, with $r=10,14,18$ even. Half this divisor defines a line bundle
whose square is trivial. The associated square-root algebra gives
a finite etale double cover on which $q$ is a square. If already a
square, use the trivial cover. A connected component of that cover
is an abelian variety after choosing an origin
\cite[Chapter I, Remark 12.3]{Milne}. Alternatively over $\mathbb C$
it is a finite-index lattice cover of a complex torus, and is
projective by pullback of an ample bundle.

If $t$ were nonconstant, this abelian variety would map rationally
to $s^2=\prod_i(t-t_i)$, a fixed smooth curve of genus
$(r-2)/2\ge4$. Composing with its Abel--Jacobi embedding gives a
rational map to its Jacobian. It extends and is a translated
homomorphism. A one-dimensional image is an elliptic translate,
which cannot be that higher-genus curve. Thus $t$ is constant.
\end{proof}

\section{Scope and exact checks}
The package \texttt{native\_pole\_ramification\_v75}, under
\nolinkurl{negative_results/computation/pole_parameters/},
checks the norm identities, the exact physical example, its point
duplication and tripling, and every native pre-pole character guard.
Finite valuation boundary checks supplement the universal proofs;
they are not used to infer them. The proof handles zero tests, trace
cancellation and zero terminal ordinates explicitly.

The ramification theorem applies to covers carrying native witnesses
or another certificate with the same valued-function-field property.
It does not cover an arbitrary arithmetically sound criterion whose
proof over $\Q$ does not extend to those geometric valuations.
The abelian theorem leaves constant slopes, individual rational
witnesses and nonabelian varieties of dimension at least two open.
None of the results disproves the required prepared realization for
integer inputs, proves a variable lower bound, or completes the
twenty-one- or nine-variable target.

\begin{thebibliography}{9}\raggedright
\bibitem{Intrinsic}\emph{An intrinsic two-parameter pole witness},
\nolinkurl{papers/pole_parameters/intrinsic_pole_20261117_v56.tex}.
\bibitem{Unit}\emph{A single norm determinant for the 25- and 13-variable constructions},
\nolinkurl{papers/finite_algebra/unit_norm_20261114_v28.tex}.
\bibitem{Seven}\emph{Seven-contact detectors and a genus-twelve actual weight curve},
\nolinkurl{papers/conditional_few_variables/seven_contact_detector_20261117_v74.tex}.
\bibitem{Even25}\emph{Even-valuation abscissa tests in the general detector},
\nolinkurl{papers/few_variables25/even_valuation25_20261117_v75.tex}.
\bibitem{Abelian}\emph{Curves and abelian families in an actual-trace square locus},
\nolinkurl{negative_results/papers/pole_parameters/curves_and_abelian_families_20261117_v69.tex}.
\bibitem{Milne}J. S. Milne, \emph{Abelian Varieties}, version 2.0,
2008, Chapter I, Theorem 3.2, Corollary 3.7 and Remark 12.3.
\url{https://www.jmilne.org/math/CourseNotes/AV.pdf}.
\bibitem{Stacks}The Stacks Project, Section 53.12, Riemann--Hurwitz,
Tag 0C1B.\quad\url{https://stacks.math.columbia.edu/tag/0C1B}.
\end{thebibliography}
\end{document}
